\chapter{无监督学习：分组与低维表示}
\label{ch:unsupervised}

聚类寻找分组，降维寻找较少的表示坐标。没有任务标签时，也要明确距离、结构假设和结果用途。
簇编号本身没有类别含义；低维图分得好看也不证明发现了真实类别。

\begin{keyidea}[先决定什么叫相似，什么值得保留]
特征选择和尺度会影响聚类；PCA保留较大方差方向。
这些目标与监督任务的预测准确性不同，需要结合用途评价。
\end{keyidea}

\section{$K$-means：分配，再更新中心}
\label{sec:unsup-kmeans}
$K$-means最小化各点到所属中心的平方欧氏距离和：
\[
 J=\sum_i\lVert\vct x_i-\vct\mu_{c_i}\rVert_2^2.
\]
常用Lloyd迭代交替进行最近中心分配与簇内均值更新。

\begin{figure}[H]
\centering\begin{tikzpicture}[x=21mm,y=11mm]
\draw[->,Quiet] (-.3,0)--(5.5,0) node[right,lab] {$x$};
\foreach \v in {0,1}{\fill[Teal] (\v,0) circle (3pt);}
\foreach \v in {4,5}{\fill[Amber] (\v,0) circle (3pt);}
\foreach \v in {0,1,4,5}{\node[lab,below] at (\v,-.1) {$\v$};}
\node[lab,above] at (0,1.4) {初始0};\node[lab,above] at (5,1.4) {初始5};
\draw[flow] (0,1.2)--(.5,.5);\draw[flow] (5,1.2)--(4.5,.5);
\node[lab,above] at (.7,.5) {新中心0.5};\node[lab,above] at (4.3,.5) {新中心4.5};
\node[lab] at (2.5,-1.1) {分配后，分别对簇内位置取均值};
\end{tikzpicture}

\caption{四个一维样本0、1、4、5，初始中心0与5；一次更新得到0.5与4.5。}
\end{figure}
\readfigure{初始分组为$\{0,1\}$、$\{4,5\}$；更新后组内平方和从2降到1。
再次分配不变，此例已经稳定。一般迭代不能保证得到全局最优，需考虑初始化和多次运行。}

平方距离偏好相应几何中的紧凑分组，异常值可显著移动均值。
空簇需要实现约定；混合尺度、非球状结构或密度差异大时，应比较其他方法与表示。
$K$-means++改善初始化选择，但不把全局最优变成保证。

\clearpage
\section{选择簇数与轮廓系数}
\label{sec:unsup-silhouette}
组内平方和一般随$K$增加而不增，因此不能直接选择其最小值来确定合理簇数。
肘部法是观察改进幅度的启发式，未必出现清晰拐点。
轮廓系数（Silhouette Coefficient）同时比较簇内和簇间距离。

对样本$i$，$a_i$是到同簇其他点的平均距离，$b_i$是到各其他簇的平均距离的最小值，
\[
 s_i=\frac{b_i-a_i}{\max(a_i,b_i)},\qquad S=\frac1n\sum_i s_i.
\]
常见取值在$[-1,1]$，较大表示在该距离定义下更接近本簇、远离最近其他簇。
单点簇通常按实现约定取0；一个簇或每点单独一簇不能按普通定义有效比较。

\begin{center}
\begin{tabular}{lrrr}
\toprule
分组$\{0,1\},\{4,5\}$ & $a_i$ & $b_i$ & $s_i$\\
\midrule
点0 & 1 & 4.5 & $7/9$\\
点1 & 1 & 3.5 & $5/7$\\
\bottomrule
\end{tabular}
\end{center}
另外两点对称，总平均为$47/63\approx0.746$。
这里用普通一维欧氏距离计算轮廓，不能把$K$-means目标中的平方距离不加说明地代入。

\quick{轮廓分数依赖距离和几何，不直接证明分组符合实际用途或真实标签。}

\clearpage
\section{DBSCAN：通过核心点扩展密度连接}
\label{sec:unsup-dbscan}
DBSCAN是基于密度的聚类，不预先指定簇数。
给定距离阈值$\varepsilon$与最少邻域点数$m$，本节约定邻域包含点自身：
$N_\varepsilon(x)=\{x_j:d(x,x_j)\leq\varepsilon\}$。

\begin{figure}[htbp]
\centering\begin{tikzpicture}[x=4cm,y=1cm]
\draw[Quiet,->] (-.15,0)--(3,0) node[right,lab] {$x$};
\draw[Teal,dashed] (.04,.85)--(.56,.85);
\draw[Teal] (.04,.73)--(.04,.97);\draw[Teal] (.56,.73)--(.56,.97);
\node[lab,above] at (.3,.95) {$0.3\pm\varepsilon$};
\foreach \x in {.1,.2,.3,2.6,2.7}{\fill[Teal] (\x,0) circle(2.7pt);}
\foreach \x in {0,.55,2.5,2.8}{\draw[Amber,thick,fill=white] (\x,0) circle(2.7pt);}
\draw[Ink,thick] (1.48,-.08)--(1.52,.08) (1.48,.08)--(1.52,-.08);
\foreach \x in {0,.1,.2,.3,.55,1.5,2.5,2.6,2.7,2.8}{\node[lab,below,rotate=55,anchor=north east] at (\x,-.07) {\x};}
\node[lab] at (1.4,-1.4) {实心：核心点\quad 空心：边界点\quad 叉号：噪声};
\end{tikzpicture}

\caption{一维例子，$\varepsilon=0.26,m=4$。点的类型由邻域计数决定。}
\par\smallskip
\readfigure{0.3的邻域为$\{0.1,0.2,0.3,0.55\}$，因此是核心点；
0.55仅有两个邻域点，是靠近核心点的边界点；1.5既非核心，也不邻接任何核心，是噪声。}
\end{figure}

\begin{center}
\begin{tabularx}{\textwidth}{P{28mm}Y}
\toprule
点的类型 & 判定与扩展\\
\midrule
核心点 & $|N_\varepsilon(x)|\geq m$，可继续扩展邻域。\\
边界点 & 不满足核心条件，但落入某核心点的邻域；可以归簇，不能继续作为扩展桥梁。\\
噪声点 & 没有归入任何密度簇的点；“噪声”不自动意味着测量错误。\\
\bottomrule
\end{tabularx}
\end{center}

先连接彼此邻接的核心点，再吸收邻近边界点，可以得到两个簇，1.5留作噪声。
边界点若同时邻接不同簇，归属可能受遍历约定影响；不能通过一串非核心边界点任意合并簇。

尺度和距离仍很关键：$\varepsilon$过小可能几乎全为噪声，过大可能合为一簇。
不同密度的簇未必能用同一阈值识别；高维距离也会变得困难。
“形状灵活”不等于对数据没有结构假设。

\section{层次聚类：链接方式决定合并距离}
\label{sec:unsup-hierarchy}
凝聚式层次聚类（Agglomerative Clustering）从每点一簇开始，
每次合并距离最近的两个簇。树状图（Dendrogram）的高度表示合并距离，
切割高度可产生不同簇数。

\begin{figure}[htbp]
\centering\begin{tikzpicture}[x=1cm,y=.8cm]
\foreach \panelshift/\mergeheight/\linkname in {0/3/单链接,7/6/全链接}{
\begin{scope}[xshift=\panelshift cm]
\draw[Quiet,->] (-.6,0)--(-.6,6.5) node[above,lab] {距离};
\foreach \y in {1,2,3,6}{\node[lab,left] at (-.6,\y) {\y};}
\draw[Teal,thick] (0,0)--(0,1)--(1.3,1)--(1.3,0);
\draw[Teal,thick] (2.6,0)--(2.6,2)--(3.9,2)--(3.9,0);
\draw[Teal,thick] (.65,1)--(.65,\mergeheight)--(3.25,\mergeheight)--(3.25,2);
\foreach \x/\label in {0/0,1.3/1,2.6/4,3.9/6}{\node[lab,below] at (\x,0) {\label};}
\node[lab,below] at (1.95,-.65) {\linkname};
\end{scope}}
\end{tikzpicture}

\caption{同样四点0、1、4、6，链接方式不同，最后合并高度不同。}
\par\smallskip
\readfigure{先在距离1处合并0与1，再在距离2处合并4与6。
随后两簇的单链接距离为3，全链接距离为6；平均链接为4.5。簇内点未移动。}
\end{figure}

单链接取跨簇点对的最小距离，容易形成链；全链接取最大距离；
平均链接取跨簇点对距离的平均。
Ward方法则以组内平方和增加为准则，常用于欧氏设置，不能混同于上述最短/最长距离。

\quick{聚类都需要选表示、尺度、距离与规则。比较结果时说明簇数、噪声处理和用途，不能只按图形好看判断。}

\begingroup\small
\noindent\textbf{参考：}Ester等，DBSCAN（1996）；ISLR第二版第12章；scikit-learn Clustering文档。
\par\endgroup

\clearpage
\section{PCA：寻找单位长度的最大方差方向}
\label{sec:unsup-pca}
主成分分析（Principal Component Analysis, PCA）先中心化输入。
第一方向最大化投影方差，约束$\lVert\vct v\rVert_2=1$；后续方向在与先前方向正交的条件下依次选择。
不加长度约束，任意放大系数就能放大方差，问题没有同样意义。

\begin{figure}[H]
\centering\begin{tikzpicture}
\begin{axis}[width=91mm,height=76mm,xmin=-2.7,xmax=2.7,ymin=-2.7,ymax=2.7,
 xlabel={$x_1$},ylabel={$x_2$},xtick={-2,0,2},ytick={-2,0,2},axis lines=middle,axis equal image]
\addplot[Quiet,dashed,domain=-2.4:2.4]{x};
\addplot[only marks,mark=*,Teal,mark size=3pt] coordinates {(2,1)(1,2)(-1,-2)(-2,-1)};
\draw[Quiet,dashed] (axis cs:2,1)--(axis cs:1.5,1.5);
\draw[Quiet,dashed] (axis cs:1,2)--(axis cs:1.5,1.5);
\draw[Quiet,dashed] (axis cs:-2,-1)--(axis cs:-1.5,-1.5);
\draw[Quiet,dashed] (axis cs:-1,-2)--(axis cs:-1.5,-1.5);
\draw[flow,very thick] (axis cs:0,0)--(axis cs:1.8,1.8) node[right,lab] {第一方向};
\draw[flow,Amber,very thick] (axis cs:0,0)--(axis cs:-.9,.9) node[left,lab] {第二方向};
\end{axis}
\end{tikzpicture}

\caption{中心为0的四点$(2,1),(1,2),(-1,-2),(-2,-1)$，第一方向为$(1,1)/\sqrt2$。}
\end{figure}
\readfigure{沿第一方向投影后，$(2,1)$的坐标为$3/\sqrt2$，重建为$(1.5,1.5)$。
虚线表示丢弃第二方向造成的误差，平方距离为0.5。第一方向描述输入变化，不直接使用类别标签。}

投影坐标为$Z=X_cV_r$，重建为$\hat X=ZV_r^\top+\vct1\vct\mu^\top$。
中心与方向均在训练内拟合。主成分符号正负可以同时翻转，不改变对应子空间。

\clearpage
\section{保留多少主成分：方差与任务一起看}
\label{sec:unsup-components}
采用分母$n-1$的样本协方差时，上页四点的矩阵为
\[
 S=\begin{bmatrix}10/3&8/3\\8/3&10/3\end{bmatrix},
\]
特征值为6与$2/3$，解释方差比例为各特征值除以总和。

\begin{figure}[H]
\centering\begin{tikzpicture}
\begin{axis}[width=119mm,height=45mm,ymin=0,ymax=1.05,xmin=.5,xmax=2.5,
 xtick={1,2},xticklabels={第一主成分,第二主成分},ytick={0,.1,.5,.9,1},
 ylabel={解释方差比例},ybar,bar width=18mm,nodes near coords,
 point meta=explicit symbolic]
\addplot[fill=Teal!65,draw=Teal] coordinates {(1,.9)[90\%](2,.1)[10\%]};
\end{axis}
\end{tikzpicture}

\caption{本例第一方向解释90\%方差，第二方向解释10\%；保留一维后仍有重建误差。}
\end{figure}
\readfigure{四条记录各损失0.5的平方距离，重建误差平方和为2。
若用“至少90\%方差”规则，本例保留一维；若要求超过90\%，则需两维。}

比例阈值是工程选择，不存在普遍正确的95\%或99\%。
较小方差方向也可能包含重要预测信息；应把PCA与后续模型一起在训练内验证。
PCA通常中心化，但是否标准化取决于单位和目标；标准化会改变所优化的方差结构。

\quick{解释方差比例描述这个表示对输入变动的保留，不等于保留同比例的“所有有用信息”。}

\clearpage
\section{核PCA与方法选择的边界}
\label{sec:unsup-kernel-pca}
核PCA（Kernel PCA）在核所对应的特征空间中进行PCA，可形成非线性表示。
训练核矩阵也要在该特征空间中心化，常写为$K_c=HKH$，
$H=I-\frac1n\vct1\vct1^\top$。

\begin{center}
\begin{tabularx}{\textwidth}{P{29mm}P{46mm}Y}
\toprule
方法 & 主要目标 & 不能直接推断\\
\midrule
$K$-means & 降低簇内平方距离。 & 得到的簇自动等于真实类别。\\
线性PCA & 用正交方向保留较大输入方差。 & 高解释方差就保留最强标签信号。\\
核PCA & 在核特征空间中寻找主成分。 & 降维后一定线性可分或可无损重建。\\
\bottomrule
\end{tabularx}
\end{center}

核与超参数决定表示；核空间的方差不能不加区别地解释为原输入单位下的方差。
新点的变换要沿用训练中心化规则；一般核表示的原空间逆映射并不自动存在。
对初学者，应先建立线性PCA与简单聚类基线，再决定是否需要非线性方法。

\subsection{自测与回看}
$\xi=0$是否只表示在SVM间隔边界？\textbf{不是。}
轮廓系数是否直接使用组内平方和？\textbf{不是。}
PCA保留90\%方差是否保证保留90\%预测能力？\textbf{不能。}
核PCA是否保证任何标签线性可分？\textbf{不能。}

\noindent\textbf{参考：}ISLR第二版第12章；MML第10章；scikit-learn Clustering与Decomposition文档。
图和数字例子独立构造。
